% 核心物理约束
\section{晶圆级交换机的核心物理约束}

晶圆上构建交换机面临五类核心物理约束。它们共享同一组物理资源，相互之间不是独立的检查和，而是紧耦合的制约网络。

\subsection{信号完整性}

每条物理链路需要若干条物理lane来实现目标带宽。
一条800\,Gbps链路若采用UCIe-32G标准（每lane 32\,Gbps），单向即需25条lane~\cite{ucie2.0-2024}，
每条lane通过一个$\mu$bump从裸片引出至interposer。
裸片的$\mu$bump总数受限于裸片面积与bump pitch——后者由所选封装工艺决定。
Dragonfly拓扑中，单个路由器裸片可连接8条以上邻居链路，
若每条链路均需25--50条lane，信号bump需求将迅速消耗裸片的总bump预算。
更关键的是，同一裸片上的$\mu$bump并非全部可用于信号互联——
部分bump必须分配给供电网络，从而在信号与电源之间形成零和竞争。

\subsection{电源完整性}

每个裸片通过$\mu$bump从interposer获取供电。
所需电源bump的数量由裸片功耗、供电电压和单bump载流能力共同决定~\cite{electrical-thermal-chiplet2024}。
裸片功耗越高，电源bump需求越大，可用于信号的bump数量越少——
两者在同一物理面积上形成零和博弈。
这一博弈的自反性在于：增加链路数量或提高每lane速率会提升裸片功耗，
进而推高电源bump需求，反过来压缩信号bump的可用面积。
性能追求与物理可行性在此处形成负反馈。

\subsection{功耗墙与热约束}

物理层互联消耗功率。UCIe标准下每条lane的PHY功耗约为5\,mW量级~\cite{ucie2.0-2024}，
计入SerDes和MAC层开销后，单条800\,Gbps链路的物理层功耗可达数瓦。
冷却方案的热移除能力受物理限制：
风冷约0.5\,W/mm$^2$、液冷约2.0\,W/mm$^2$、浸没冷却约5.0\,W/mm$^2$、微流道约10.0\,W/mm$^2$。
更大规模的拓扑和更密集的功耗分布将迅速逼近给定冷却方案的热上限。

然而功耗上限仅是热约束的一个维度。
更棘手的是功耗的空间不均匀性——热量在硅interposer上横向传导，
裸片间热串扰使得部分裸片的局部温度显著高于晶圆平均温度~\cite{mfit2025}。
由此产生的温度梯度引发差异化的热膨胀，
在$\mu$bump与interposer界面产生剪切应力，
当应力超过疲劳极限时导致bump断裂或界面脱层——即翘曲失效。
因此，热约束同时包含总功耗上限和空间温差上限两个维度，
后者对功耗分布的均匀性施加了额外限制。

\subsection{布线密度}

除bump资源外，interposer金属层的走线容量构成另一项几何约束。
裸片布局确定后，每条interposer grid边上可容纳的走线数量受限于可用金属层数和每层走线密度。
典型晶圆级interposer提供4--6层金属用于信号走线，每毫米约承载100--200条走线。
若多条incident链路的Manhattan走线路径汇聚于同一grid边，
该边的走线需求可能超出物理容量。
在早期设计阶段，完整的布局布线不可行，
因此需要一种粗粒度但能够捕捉走线拥塞趋势的容量估计方法。

\subsection{耦合的本质}

上述五类约束构成一个闭环的耦合网络。
信号bump数量取决于电源bump需求（电源 $\to$ 信号），
电源bump需求取决于裸片功耗（热 $\to$ 电源），
裸片功耗取决于链路数量和速率（信号 $\to$ 热），
链路负载的空间分布决定温差分布（流量 $\to$ 翘曲），
布线容量约束了bump所能连接的链路数量（布线 $\to$ 信号）。
任何单一维度的独立优化都可能被另一维度的耦合效应所推翻，
使得"先设计后验证"的序贯流程在晶圆级交换机场景下不再成立。
